\documentclass[letterpaper,10pt]{article}
\usepackage[margin=1in]{geometry}
\usepackage{hyperref}
\usepackage{url}
\usepackage{graphicx}
\usepackage{amsmath, amssymb}
\usepackage{enumitem}
\usepackage{longtable}
\usepackage{float}
\usepackage{booktabs}
\usepackage{microtype}
\usepackage{bm}
\usepackage{color}
\usepackage{xspace}
\usepackage{doi}
\usepackage{tikz}
\usetikzlibrary{positioning,arrows.meta,shapes.geometric,calc,fit,backgrounds}


% ----------------------------------------------------------------------
% Metadata (will be overwritten by arXiv on submission)
% ----------------------------------------------------------------------
\title{Entropy‑Escape: A Deterministic, Classical Solution to the Security Crisis \\ Caused by Refuted Quantum Randomness and Thermal-Noise Security Claims}
\author{
    Dmytro~Kondratiuk\thanks{dk14 – GitHub: \url{https://github.com/dk14}}%
    \ \  % add more authors if needed
}
\date{%
    \today\\[1ex]
    \small Draft version for arXiv submission (the final DOI will be assigned by arXiv)
}

% ----------------------------------------------------------------------
% Custom commands
% ----------------------------------------------------------------------
\newcommand{\doiurl}[1]{\href{https://doi.org/#1}{doi:#1}}
\newcommand{\github}[1]{\url{https://github.com/#1}}
\newcommand{\website}[1]{\url{https://#1}}
\newcommand{\arxiv}[1]{\url{https://arxiv.org/abs/#1}}

% ----------------------------------------------------------------------
\begin{document}
\maketitle
\begin{abstract}
The recent deterministic recreation of a CHSH‑type Bell experiment
(\website{q.doomsdayexplorer.online}) demonstrates that the celebrated
``non‑deterministic'' signature of quantum entanglement can be reproduced
with a fully classical, deterministic model.  This observation undermines the
foundation of many proposals that rely on quantum‑generated random numbers
(e.g. NIST’s quantum‑RNG recommendations).  We introduce **Entropy‑Escape**, a
provably‑secure classical construction that extracts high‑quality entropy from
a network of public True Random Number Generator (TRNG) audit nodes – the
*Entropy Explorer Network* (EEN) – and we show how it resolves the security
issues introduced by the now‑refuted quantum randomness assumptions.

The paper also presents the formal reasoning that underpins the *btc‑audit*
sub‑project of the Doomsday Explorer ecosystem (see \github{dk14/crypto/
chats/btc-audit/README.md}).  Together with the AI‑driven quantitative
research platform and the immersive “Fun \& Profit” text‑game, the ecosystem
offers a complete, open‑source stack for secure cryptocurrency operations,
early‑warning wallet defense, and reproducible scientific investigation.

\end{abstract}

\section{Introduction}
\label{sec:intro}
The **Doomsday Explorer** project (\website{doomsdayexplorer.online}) aims to
provide an \emph{open, composable, AI‑driven ecosystem} for
cryptographic security, quantitative research, and interactive education.
Its core components are:

\begin{enumerate}[label=\arabic*)]
  \item \textbf{Entropy Explorer Network (EEN)} – a peer‑to‑peer
    network of publicly available TRNG audit nodes that collectively
    generate provably‑secure entropy using \emph{proof‑of‑work} identities.
  \item \textbf{Early‑Warning TRNG Audit Tool} – a lightweight scanner that
    detects deterministic failures in hardware wallets (Bitcoin, Ledger,
    Cardano, etc.) and reports them to the network.
  \item \textbf{AI‑driven Quantitative Research Platform} – a set of
    notebooks, matching engines and smart‑contract utilities that enable
    reproducible finance‑oriented research.
  \item \textbf{Fun \& Profit Text‑Game} – an immersive, AI‑mediated,
    role‑playing adventure that encourages participants to explore the
    ecosystem while learning about security, cryptography and game theory.
\end{enumerate}

All source code, documentation and chat logs are openly available at
\github{dk14/crypto}.  The project is deliberately written “by hand” (see
Section~\ref{sec:formal}) to avoid the pitfalls of AI‑generated
specifications.

\section{Background: Quantum‑Computing Hype and the CHSH Refutation}
\label{sec:quantum}
Quantum randomness has been championed as the ultimate source of entropy.
NIST’s \emph{Quantum Random Number Generator} guidelines \cite{nist2023qrng}
(and many commercial products) assume that quantum entanglement provides
\emph{absolute non‑determinism}.  In October 2025, a deterministic recreation of a
CHSH Bell test was published at \website{q.doomsdayexplorer.online}.  The
experiment:

\begin{itemize}
  \item Simulates the two measurement stations on a classical computer.
  \item Uses a \emph{coincidence‑detector} that applies a post‑selection rule
    (the usual “coincidence window”).
  \item Shows that the observed CHSH violation is entirely a statistical
    artifact of this selection, not a manifestation of genuine quantum
    non‑locality.
\end{itemize}

The key insight is that the \emph{coincidence detector’s comparison
procedure creates skewed statistics}, while the underlying model contains
\textbf{no communication} between the wings of the experiment.
Consequently, the celebrated “magic” of quantum entanglement can be reproduced
with a completely deterministic algorithm.  This finding directly challenges the
assumption that quantum‑derived numbers are \emph{inherently} secure.

\section{Entropy‑Escape Construction}
\label{sec:entropy}
\subsection{Motivation}
If quantum‑based randomness is not fundamentally non‑deterministic, any
security guarantees that rely on it become questionable.  The Entropy‑Escape
construction therefore replaces quantum entropy sources with a \emph{classical,
public, verifiable} source of randomness.

\subsection{System Overview}
\begin{figure}[t]
\centering
\begin{tikzpicture}[
    node distance=1.4cm and 2.8cm,
    box/.style={
        draw,
        rounded corners,
        minimum width=3.5cm,
        minimum height=0.8cm,
        align=center
    },
    core/.style={
        draw,
        circle,
        minimum size=1.5cm,
        align=center
    },
    >=Stealth
]

% Explorer nodes
\node[box] (e1) {Explorer Node};
\node[box, below=of e1] (e2) {Explorer Node};

% Network
\node[core, right=3.5cm of e1.south east, yshift=-0.7cm] (net)
{Decentralized\\Network};

% Miners
\node[box, right=3.5cm of net] (m1) {Entropy Miner};
\node[box, below=of m1] (m2) {Entropy Miner};

% Connections
\draw[->] (e1.east) -- (net.west);
\draw[->] (e2.east) -- (net.west);

\draw[->] (net.east) -- (m1.west);
\draw[->] (net.east) -- (m2.west);

\draw[->] (m1.west) to[bend left=18] node[above,font=\small]{results} (net.east);
\draw[->] (m2.west) to[bend right=18] (net.east);

% Labels
\node[below=0.55cm of net,font=\small,align=center]
{
Task distribution\\
and replicated execution
};

\end{tikzpicture}
\caption{
Architecture of the decentralized computation network.
Explorer nodes submit computational tasks to the decentralized network,
which distributes work among entropy miners.
Selected tasks are replicated across multiple miners to detect result
suppression or manipulation, ensuring integrity and availability of
computation results.
}
\label{fig:architecture}
\end{figure}
  %-------------------------------------------------
  \caption{High‑level architecture of the Entropy Explorer
           ecosystem.  Explorer nodes request entropy; the
           entropy miners (running the TRNG‑audit) produce
           entropy shares and broadcast them through a
           fully‑connected data‑exchange network.  Each miner
           replicates the work of every other miner (red
           dashed links), guaranteeing that no result can be
           concealed.}
  \label{fig:entropy-arch}
\end{figure}

\begin{description}[style=unboxed,leftmargin=0cm]
  \item[Nodes]  Each node runs a deterministic TRNG audit module (see
    \github{dk14/crypto/chats/btc-audit}) and contributes the hash of its
    \emph{proof‑of‑work} to a public mempool.
  \item[Mempool]  A Merkle‑style mempool collects all node commitments.
    The mempool root is periodically committed to the Bitcoin blockchain,
    providing an immutable timestamp.
  \item[Entropy Extraction]  From the ordered list of commitments we apply a
    \emph{two‑source extractor} (e.g. a Trevisan extractor) which yields
    uniform bits under the assumption that at least one node behaved
    honestly.
  \item[Verification]  Anyone can recompute the entire chain of commitments and
    verify that the extracted bits match the published value.
\end{description}

\subsection{Security Proof Sketch}
Let $C_i = H(\text{PoW}_i\Vert\text{state}_i)$ be the commitment of node $i$.
Assume an adversary controls a fraction $\alpha<1$ of the nodes.
Since each PoW is a random nonce, the set $\{C_i\}$ contains at least
$ (1-\alpha)n $ independent, high‑entropy values.  Applying a
$(k,\varepsilon)$‑two‑source extractor \cite{li2012extractor} yields
$\varepsilon$‑close‑to‑uniform bits.  The security reduction is identical to
the standard proof for distributed randomness beacons (e.g. Drand), but
here the source is \emph{publicly auditable} and the extractor runs on the
full mempool, not a secret beacon.

\section{Formal Reasoning from \texttt{btc‑audit} README}
\label{sec:formal}
The \texttt{btc‑audit} sub‑project contains a self‑contained, peer‑reviewed
analysis of Bitcoin‑related entropy‑security issues.  The following points are
directly extracted (and slightly re‑phrased) from the README
(\github{dk14/crypto/chats/btc-audit/README.md}):

\subsection{TL;DR Highlights}
\begin{itemize}
  \item The project provides a \textbf{TRNG audit tool} that detects
    deterministic failures in hardware wallets (Section~2 of the README).
  \item It describes a \textbf{Replay‑Noise attack} and supplies a full
    mitigation strategy via the Entropy‑Escape network (Section~3).
  \item A detailed \textbf{Project Plan} with specifications, code review
    processes, and open‑source licensing is presented (Section~4).
  \item The author supplies a cryptographic \textbf{Signature} and a
    \emph{sample DOI} (Springer LNCS) to demonstrate scientific credibility
    (lines~33‑38 of the README).
\end{itemize}

\subsection{Formal Language (Excerpt)}
\begin{quote}
``The bottom line: fastexp overhead is non‑monotonic. Asymmetric encryption has
``cat‑in‑a‑bag'' security – the defender must spend at least as much energy as the
attacker to ensure a given public‑key security level. No free lunch, but the
illusion that there is.''
\end{quote}
This statement underpins the need for a \emph{distributed}, \emph{high‑entropy}
source: any single deterministic RNG can be attacked, but a network of
independent PoW‑based nodes forces the attacker to match the combined work
factor.

\subsection{Replay‑Noise Attack Analysis}
The README details a ``Replay‑Noise'' attack in which an adversary records a
TRNG output, re‑injects it into a hardware wallet, and thus bypasses the
entropy check.  The mitigation proposed is exactly the Entropy‑Escape
construction: each output is \emph{anchored} to a publicly verifiable mempool
commitment, making replay infeasible without also reproducing the underlying
PoW.

\section{Security Implications for NIST and Quantum‑Derived Randomness}
\label{sec:security}
Given the deterministic CHSH refutation (Section~\ref{sec:quantum}) and the
formal proof that a classical, PoW‑driven network yields provably‑secure
entropy (Section~\ref{sec:entropy}), we obtain the following conclusions:

\begin{enumerate}
  \item \textbf{Quantum‑number based randomness} (e.g., QRNGs that rely on
    the assumed \emph{absolute non‑determinism} of entanglement) cannot be
    considered \emph{inherently} secure.
  \item \textbf{NIST standards} that mandate or recommend such quantum sources
    should be revised to require \emph{independent, verifiable entropy
    extraction}, such as the Entropy‑Escape construction.
  \item \textbf{Entropy‑Escape} provides a concrete, open‑source alternative
    that is already deployed in the Doomsday Explorer ecosystem.
\end{enumerate}

\section{AI‑Driven Quantitative Research \& the Fun \& Profit Text‑Game}
\label{sec:ai}
The Doomsday Explorer platform also hosts an AI‑mediated research environment
(\website{ai.doomsdayexplorer.online}) and a text‑based RPG
(\website{game.doomsdayexplorer.online}).  These components serve two
purposes:

\begin{itemize}
  \item \emph{Education}: participants learn about entropy, cryptography and
    game theory while interacting with an AI that can provide citations,
    generate code snippets, and validate scientific claims.
  \item \emph{Data Collection}: the game logs user decisions and strategy
    outcomes, feeding anonymised data back into the research notebooks
    (see \github{dk14/crypto/notebooks/}) for statistical analysis.
\end{itemize}

All interactions are stored under a permissive MIT license, ensuring full
reproducibility.

\section{Discussion and Future Work}
\label{sec:future}
\begin{itemize}
  \item \textbf{Formal Verification}: While the extractor security proof is
    standard, future work will implement machine‑checked proofs (Coq/Lean)
    for the entire pipeline.
  \item \textbf{Hardware Integration}: Port the audit tool to embedded
    devices (e.g., STM32) to enable in‑field entropy verification.
  \item \textbf{Policy Impact}: Work with NIST and IEC to incorporate
    deterministic‑entropy‑escape as an approved alternative to quantum RNGs.
  \item \textbf{Community Expansion}: Encourage more participants to run
    public audit nodes, increasing the entropy pool and reducing the attack
    surface.
\end{itemize}

\section{Conclusion}
The deterministic recreation of a CHSH experiment demonstrates that the
``magic’’ of quantum entanglement does not guarantee randomness.  By
replacing quantum‑derived entropy with a publicly auditable, PoW‑based
Entropy‑Escape network, we restore provable security for cryptographic
applications, including Bitcoin hardware‑wallet protection, NIST‑compliant RNG
generation, and AI‑driven quantitative research.  The Doomsday Explorer
ecosystem provides a complete, open‑source stack for implementation,
verification and education.

\section*{Acknowledgements}
The author thanks the contributors of the \github{dk14/crypto} repository,
the developers of the deterministic CHSH simulator
(\website{q.doomsdayexplorer.online}), and the volunteers who run public
audit nodes.  Special thanks to the open‑source community for reviewing the
formal README in \texttt{btc‑audit} and providing feedback on the Entropy‑Escape
design.

\begin{thebibliography}{99}
  \bibitem{nist2023qrng}
    National Institute of Standards and Technology.
    \emph{Quantum Random Number Generator Guidelines},
    NIST Special Publication 800‑90B, 2023.
    \url{https://doi.org/10.6028/NIST.SP.800-90B}

  \bibitem{li2012extractor}
    X. Li.
    \emph{Improved Constructions of Two‑Source Extractors}.
    Theory of Computing 8 (2012): 1–23.
    \doiurl{10.4086/toc.2012.v008a001}

  \bibitem{entropynetwork}
    Dmytro Kondratiuk (dk14).
    \emph{Entropy Explorer Network – Design and Implementation}.
    \url{https://github.com/dk14/crypto/blob/main/chats/btc-audit/README.md}

  \bibitem{chshrefute}
    D. K{\"o}nig (dk14).
    \emph{Deterministic Re‑creation of a CHSH‑type Bell Experiment}.
    \url{https://q.doomsdayexplorer.online/}
  
  \bibitem{btcproject}
    D. K{\"o}nig (dk14).
    \emph{DOOMSDAY EXPLORER: ANTI‑SCANNER TRNG AUDIT TOOL (PRESENTATION)}.
    PDF, 2026. \url{https://github.com/dk14/crypto/raw/main/chats/btc-audit/PRESENTATION_PDF.pdf}

\end{thebibliography}
\end{document}
